\documentclass[a4paper,10pt]{article}

\usepackage{fcds}
\usepackage{graphicx}
\usepackage{float}
\usepackage[labelsep=period]{caption}
\usepackage{enumitem}
\usepackage[utf8]{inputenc}
\usepackage{titlesec}
\titlelabel{\thetitle.\quad}
\usepackage[top=5.5cm, bottom=5cm, left=4cm, right=4cm]{geometry}
\graphicspath{ {./pics/} }
\newcommand{\ax}{a\hspace{-0.45mm}{\scriptsize $\!_{\mbox{\char 62}}$}
\hspace{-0.31mm}}
\date{}
\providecommand{\keywords}[1]{\textbf{{Keywords---}} #1}
%\renewcommand{\thesubsection}{\arabic{subsection}}

\usepackage[colorlinks=true]{hyperref}
\setcounter{tocdepth}{3}
	

\title{\textbf{Using \TeX \space markup language for 3D and 2D geological plotting by packages \{pgfplots\}, \{tikz\} and \{tikz-3dplot\} }}

\author{ $\,\,\,\,\,\,\,$ Vasiliy Lemenkov, Polina Lemenkova \address{Geoprojectsurvey LLC, Moscow, Russia and Ocean University of China, College of Marine Geo-sciences, Qingdao, China, \{wasiliy.lemenkov, pauline.lemenkova\}@gmail.com}}

\begin{document}

\maketitle

\begin{abstract}
The paper presents technical application of \TeX \space high-level, descriptive markup language for processing geological dataset from soil laboratory. Geotechnical measurements included equivalent soil cohesion, absolute and absolute deformation index, soil compressibility coefficient by time of immersion depth, exposure time to compressive strength to samples and physical and mechanical properties (humidity, density). Dataset was received from laboratory based experimental tests of the physical and mechanical properties of soils. Data were converted to csv table and processed by \LaTeX. Methodology is base don \LaTeX \space packages: \{tikz\}, \{tikz-3dplot\}, \{tikzpicture\}, \{pgfplot\}, \{filecontetns\}, \{spy\} for 3D plotting showing correlation in variables and descriptive statistical analysis based on the data array processing. Results demonstrated \LaTeX \space scripts and graphics: 2D and 3D scatterplots, ternaries, bar charts, boxplots, zooming techniques  detailing fragment of the plot, flowchart. Research novelty consists in technical approach of \TeX \space  language application for geological data processing and graphical visualization. Engineering graphics by \TeX \space was demonstrated with screenshots of the codes used for plotting. 
\end{abstract}

\begin{keywords}
\LaTeX, \TeX, markup language, plotting, data visualization, data analysis, geotechnical modelling, soil compression, 3D modelling, 2D modelling
\end{keywords}


\section{Introduction}

The question of which software to use for geological data processing received from series of experiment in a soil laboratory is not an unambiguous one. 

On the one hand, there is traditional software where one can build graphs, tables and carry out primitive data processing. In case quick visualization of a small amount of data is required, basic standard programs (e.g. SPSS, MS Excel) can be applied to visualize correlation and plotting data distribution. It could be line of dot graphs showing variations of category X on Y. However, the graphical view and design representation of such graphs remains questionable. For more advanced cases when complex data analysis is required, high-level programming languages that can be used. 

On the other hand, if we are dealing with a large heterogeneous data array and the output requires modeling of different categories, plotting facetted plots, as well as high-quality printer-quality graphics, using advantage of the statistical libraries of the programming languages is evident (e.g. R, Python, Octave, Matlab), since they combine both mathematical powerful instrument for data analysis and  modelling and high quality, data visualization. 

\LaTeX, high-level, descriptive markup language combines graphical quality of the data visualization similar to the best statistical packages of R pf Python, with simplicity and straightforward approach to reading tables, description statistical data visualization and plotting. Hence, it can be effectively used for quick yet effective data visualization.  This paper presents technical application of \TeX \space for processing geological dataset received from the experimental series of soil laboratory.

\section{Methodology}

\subsection{\LaTeX \space packages}

Prerequisites for the methodology is \LaTeX \space Installation. Main \LaTeX \space package used is \{pgfplot\} which  is a free (open source) software. It works under various OS platforms: Mac OS, Linux, and Windows. Thus, for MacOS users MacTeX release exist. To work with LaTeX, the latest version of TeXShop Editor was installed and \LaTeX \space compiler was used. It contains necessary packages as well as additional required packages: https://www.latex-project.org.  The abbreviation 'pgf means portable graphics format'. In this study methodology includes applying \LaTeX \space packages:

\begin{itemize}
	\item  \{tikz\}
	\item \{tikz-3dplot\}
	\item \{tikzpicture\}
	\item \{pgfplot\}
	\item \{filecontetns\}, 
	\item \{spy\} 
\end{itemize}
3D visualization showing correlation between variables and descriptive statistical analysis based on the processing of a data array. 

\subsection{Reading Data}

Dataset in this study was received from experimental testing of the physical and mechanical properties of soils in laboratory. Data were converted to csv table and processed by \LaTeX. Reading data from a .csv file was done using package \{filecontents\}. 

\begin{figure}[H]
    \centering
    \includegraphics[width=0.7\textwidth]{Fig-1.jpg}
    \caption{Variability of soil properties during compression strength testing. Visualized by \LaTeX. Code for plotting graph and graphical output. Source: authors.}
    \label{fig:1}
\end{figure}

By this package the data can be read in to the \LaTeX \space environment and used to plot graphs. In this particular case, we analyzed changes incompressibility of soils under deformation conditions. 

Suppose we have a small table with three categories (conditionally, XYZ). In this simple case, we can directly add these three variable values as coordinates. In this case, they will not necessarily be geographical, but simply reflect the dependence of category values on each other: soil properties, deformation, compressibility, heave of the frozen soils depending on the temperature, and other physical and mechanical parameters of souls: pressure, humidity, immersion time, density, etc.

Figure~\ref{fig:1} presents variability of soil properties during compression strength testing. 

\subsection{Geological measurements}

Following variables were measured during the set of experiments and visualized by \LaTeX \space based plotting:
\begin{enumerate}[leftmargin=1cm]
  \item equivalent soil cohesion,
  \item absolute and absolute deformation indeces
  \item time of immersion depth
  \item compressive strength to samples
  \item physical and mechanical properties (humidity, density)
\end{enumerate}

In case of a complicated table with rows of variables and varying indicators, we used \{filecontetns\} package in the preamble of the document description environment. The \{filecontetns\} package allows \LaTeX\space to read data from long tables with many variables. In this case, we directly indicated the table in the preamble of the code (see the picture below, Figure~\ref{fig:1}). The environment between 'begin \{filecontetns\}' and 'end \{filecontetns\}' indicates which file we are working with (.csv), what variables are in, and which of them should be used for plotting. In this example, from the whole long table, we selected only 3 parameters for visualization. 

\section{Results}

\subsection{Summary}
The results are shown in 10 Figures demonstrating \LaTeX \space script and graphical output: 2D and 3D scatterplots, ternaries, bar charts, boxplots ( probability distribution), using techniques of zooming in a plot detailing fragment of the graph using magnifying glass function, and flowchart by library \{tikz\}. The screenshots show codes and graphical results: 2D and 3D scatterplots, ternaries, bar charts, boxplots ( probability distribution), using techniques of zooming in a plot detailing fragment of the graph using magnifying glass function, and flowchart by library \{tikz\}.

2D (Figure~\ref{fig:1}, plot left, Figure~\ref{fig:5}) and 3D graphs (Figure~\ref{fig:1}, plot right, Figure~\ref{fig:2}, Figure~\ref{fig:3} and Figure~\ref{fig:4}) for testing soil samples were plotted against compression strength with determination of the absolute strain index and compressibility coefficient. We have 6 variables and a time category (maximum - 600 hours for testing). In the description of the graph environment, we indicates parameters of the graph in the commands shown in Figure~\ref{fig:1} and plot the graph.

\begin{figure}[H]
    \centering
    	\includegraphics[width=0.7\textwidth]{Fig-2.jpg}
    \caption{Code for plotting 3D graph of variability of equivalent cohesion of soil in \LaTeX. Source: authors.}
    \label{fig:2}
\end{figure}

\subsection{3D scatterplot of the correlation of the equivalent cohesion of soils}
Figure~\ref{fig:2} shows workflow for plotting 3D scatterplot of the correlation between equivalent cohesion of soils, time of experiment load pressure and immersion depth of the external loads. A code with explanations (with comments highlighted in red) is describing stepwise which commands were used and parameter selected to plot two correlation graphs. Figure~\ref{fig:2} presents code and graphical plot showing 3D graph of variability of the equivalent cohesion of soil in \LaTeX. 

\begin{figure}[H]
    \centering
    	\includegraphics[width=0.7\textwidth]{Fig-3.jpg}
    \caption{Code for plotting zoomed enlarged fragment of a 3D graph of the equivalent soil cohesion with immersion depth changing in time and graphical result. Source: authors.}
    \label{fig:3}
\end{figure}

The output shows the dependences of the soil parameters against external factors (in this case, the depth of immersion of samples over time shows us resistivity of the soils), as presented below. In this case, two 3D graphics from a series of tests were placed as a facetted plot. Following requirements for adjusting plots are demonstrated on the Figure~\ref{fig:2}: rotated axes; rotated axes annotations placed below the axes; small-sized two plotted located sidewise; plotted coordinate ticks; displayed grid; title annotated on top. The source code and the output are visualized on Figure~\ref{fig:2}.

\subsection{Selective zooming of the enlarged fragment}
Figure~\ref{fig:3} presents a \LaTeX \space code for plotting zoomed enlarged fragment of a 3D graph and resulting plot. Creating an enlarged fragment of the graph has been done by \{spy\} package and magnifier command. Using magnifying glass function was applied to enlarge a fragment of the graph for a more detailed analysis, as shown on Figure~\ref{fig:3} which presents source code above and output image below. 

\subsection{Facetted plotting}

\begin{figure}[H]
    \centering
    	\includegraphics[width=0.8\textwidth]{Fig-4.jpg}
    \caption{Facetted plot of the several graphics on experimental sets: equivalent soil cohesion, immersion depth and time. Source: authors.}
    \label{fig:4}
\end{figure}

The principle of the approach is the same for a time series of such graphs Figure~\ref{fig:4} showing facetted plot of six graphs for analyzing the trend of soil behavior. Here we only placed a series of graphs one after another declaring a \{tikzpicture\} environment from the \{tikz\} package. The space between the plots is indicated by \{axis\} and 'hskip' command that inserts a non-breaking space of specified width between the graphical objects.

Figure~\ref{fig:4} presents a Facetted plot of the several graphics on experimental sets: equivalent soil cohesion, immersion depth and time. Figure~\ref{fig:5} presents a scatter 2D plot of 8-hour equivalent soil cohesion of soil samples against moisture (left) and density (right). The distribution series of the soils parameters are shown on the diagrams of 1D data distribution and a scatter plot in the \{pgfplot\} \LaTeX \space package. Main tasks of the experiments performed in soil laboratory of engineering and geological surveys include assessment of the soil properties, structure and behavior under varying conditions (pressure). 

Measurements of these parameters are performed under various environmental conditions for the soils taken at different depths of the wells, as well as factors affecting change in the soil characteristics with aim to model and numerically assess soil behavior for the construction purposes. Soil samples as study objects are heterogeneous multicomponent open dynamic natural systems. Therefore, modeling and visualizing soil compression for road construction purposes requires use of numerical methods of mathematical-statistical analysis, expressed also in a graphical form.

\subsection{Clustering data}

Another important task is determining general nature of the data distribution to assess the degree of its homogeneity. Statistical homogeneity of the measured values is characterized by the range of variation (scattering) of data. In this case we test soil parameters. The mismatch between the values of the specific variables for statistical units in clusters is assessed as degree of data dispersion.

This and other tasks of basic statistical data processing and visualization can be solved using \LaTeX \space functional. It should be however noted that complex statistical data analysis including sophisticated methods and mathematical approaches can only be solved using specialized tools, e. g. StatsModel library of Python or powerful functionality of R language. \LaTeX \space functionality is more straightforward and can be effectively used for quick data processing and visualization which enables to avoid learning complex coding syntax of Python of R, yet provides excellent graphical data visualization.

Consider a specific example. Suppose we have a series of measurements showing values of the equivalent cohesion of the soil samples, which vary with different loads. We need to visualize a series of tests and evaluate the normality of the data distribution. Visualizing series of tests and evaluating the normality of the distribution can be done by \LaTeX. Using pgf functional, the plot of values by clusters was visualized (Figure~\ref{fig:5}). Figure~\ref{fig:5}. In this case, the presence of classes of objects that are explicitly defined in the code will be critical: table (meta = class). Following statistical values from the table were entered: median, the minimum and maximum values of the data (the lowest and highest values of the array), lower and upper quartiles (values of the array corresponding to 25\% for lower quantile and 75\% for upper quantile), frequencies on the cumulative graph \cite{Kaputin}, Figure~\ref{fig:5}.

\begin{figure}[H]
    \centering
    	\includegraphics[width=0.7\textwidth]{Fig-5.jpg}
    \caption{Scatter 2D plot of 8-hour equivalent soil cohesion of soil samples against moisture (left) and density (right) from data clusters. Source: authors.}
    \label{fig:5}
\end{figure}

\subsection{Statistical boxplots}

The most common case of solving problems in the engineering geology is a direct task: according to the given characteristics of soils and their reactions to external impact factors we can study behavior of the soils samples in space and time using compression of the samples. Measured variables include deformation, creep, ductile limit for for cohesive soils, heave, cohesion, etc. There are also reverse and inverse geological problems, for example, assessing load pressure on samples by measured properties of soils \cite{Gorokhovskiy}. However, in all the cases the basic essence of the problems is the same: interaction between soil and external pressure, correlations, and reaction of the objects in the dynamic system.

\begin{figure}[H]
    \centering
    	\includegraphics[width=0.7\textwidth]{Fig-6.jpg}
    \caption{Plotting whisker boxplot: a diagram of 1D distribution of statistical probabilities of an 8-hour equivalent cohesion of the soil samples: a series of 15 measurements. Source: authors.}
    \label{fig:6}
\end{figure}

Estimation of the data distribution in space is common task in statistical data analysis \cite{Gorokhovsky}. In this case, a scatterplot or boxplot are often plotted based on a series of measured data (Figure~\ref{fig:5}). In a first approximation, the graph has an ellipse form. In various tasks, the parameters of the ellipse covering points cloud are required for plotting. In other cases, displaying boxplot of descriptive statistical data is effective. An ideal boxplot represent statistical variables (median, quartiles, mean) based on data distribution \cite{Porotov}.

Figure~\ref{fig:6} presents a whisker boxplot showing a diagram of 1D distribution of statistical probabilities of an 8-hour equivalent cohesion of the soil samples: a series of 15 measurements. As a general case, the points of cohesion are denser in a center of the box in a boxplot with their density decreased towards the edges and having outliers, which indicates normal distribution of the data. In other cases, distribution of the data samples can be graphically displayed in a straight line as correlation of the measurements against variable. In this case, boxplot shows distribution of the equivalent soil cohesion under varied compression strength.

\subsection{Ternary diagrams by \LaTeX}

Among the problems of modeling in geological survey, there is one interesting and less usual comparing to the linear plots: ternary diagrams. Ternary plot looks like elegant, logically constructed triangle showing triple correlation between the data variables. Since its visualization is not a classic graph showing 2D dependence of x on x or 3D, plotting ternaries \LaTeX \space requires some explanation. 

Omitting the solutions possible in R programming which require more advance skills in coding, we focus on more straightforward functionality of \LaTeX \space \{pgfplot\} and its \{ternary\} library. Ternary diagram allows to visualize variations in the lithological composition of soils or rocks of the given territory or to trace their evolution during geological time in one area through comparing ternaries of the rocks of the same age in a cross-section.

Ternary diagrams are very convenient for visualizing lithological composition of rocks and soil samples. Using ternary technique enables to visualize in form of a cloud of points any three-component system which elements which in total compose 100\% of the system's structure. Hence, each of the peaks corresponds to 100\% of one of the three components, and the side opposite it corresponds to the zero content of the same component. Figure~\ref{fig:7} and Figure~\ref{fig:8} present two ternary diagrams showing granulometric composition of the soil with Figure~\ref{fig:8} using color-plotted ternary. 

Plotting ternary diagram showing lithological composition of rocks or soils by \LaTeX \{pgfplot\} can be done by special libraries: pgfplotslibrary \{ternary\} and tikzlibrary \{pgfplots.ternary\}. In \LaTeX \space environment, ternary diagram visualizes a three-component rock system in such a way that the sum of all these components gives 100\%. Therefore, it is logical that ternary diagram visually represents a system located on the triangular axes. 

Hence, \{ternaryaxis\} environment of \LaTeX \space works with relative coordinates: each data point consists of three components x, y, z. Their total amounts to 100\%. So, the while data on one axis stretches in any of these three components the other two will automatically decrease, but in total the whole will be 100\%. The convenience in ternary approach lies in its specifics: if one chooses any absolute values within the data range, \LaTeX \space will still recalculate them into a system of barycentric (trilinear) coordinates. Therefore, in a standard configuration, we have x, y, z placed within the range (0,1).

Plotting ternary diagram was done by executing the commands demonstrated in the code Figure~\ref{fig:7} and Figure~\ref{fig:8}. Explanations of the specific commands is highlighted in red. The printscreens on Figure~\ref{fig:7} and Figure~\ref{fig:8} shows in detail (with comment explanations in red color) the techniques of ternary diagram. 

\begin{figure}[H]
    \centering
    	\includegraphics[width=0.7\textwidth]{Fig-7.jpg}
    \caption{Ternary diagram of the granulometric composition. Source: authors.}
    \label{fig:7}
\end{figure}

At the next step, ternary was plotted in color Figure~\ref{fig:8} using the same data with output shaded by colors selected by diverse color palette.  the description and codes are presented here: http://latexcolor.com/ (accessed 16.04.2020) enabling to select color both by name, by Hex RGB Triplet, or by defining individual colors through \LaTeX \space 'definecolor' declaration.
For example, used color 'Dark Mandarin' was done using following example. In the document preamble the RGB triplet was defined explicitly, and the color was then used directly in the work (here: fill = darktangerine).

\begin{figure}[H]
    \centering
    	\includegraphics[width=0.7\textwidth]{Fig-8.jpg}
    \caption{Ternary diagram showing granulometric composition of the soil. Source: authors.}
    \label{fig:8}
\end{figure}

Since rocks usually consist of many components, visual arrangement of the granulometric ternaries was divided into three groups according to dialogical and genetic characteristics. For example, in case of terrigenous rocks, all fractions of sand were combined in one group (bottom right). Second group included all fractions of silt. Third group included pelite and clay material (clay and clay rocks). Depending on goals, data could also be grouped according using other properties, e.g.: carbonate rocks, calcite, dolomite and insoluble minerals can be grouped as independent clusters. \LaTeX \space \{ternary\} library can also apply contours and smooth lines function.

\subsection{Bar charts on geological variables}

Figure~\ref{fig:9} presents a 2D bar charts showing variables measured against time and compressive strength.
Let's look at another example: we have a task to build a data distribution diagram (Box Plot, or as it is often called a 'whisker plot').

\begin{figure}[H]
    \centering
    	\includegraphics[width=0.7\textwidth]{Fig-9.jpg}
    \caption{2D bar charts showing variables measured against time and compressive strength. Source: authors.}
    \label{fig:9}
\end{figure}

To do this, \LaTeX \space function 'boxplot' was used. According to the description of the code, the graph is displayed automatically (see illustrations below). The source code and the result of the statistical data distribution diagram are shown in Fig. 5. Important technical point: \LaTeX \space can read data from the files in .dat or .csv (comma separated value) formats. In the latter case, when importing data from Excel tables, one need to check that distances between the columns are comma separated. Using available tabular data on measurements of equivalent soil cohesion (in this case, tests were carried out using equipment developed by KrioLab). The data were filled directly into the body of code: upper and lower quartiles, upper and lower 'whiskers', extreme values for the entire data sample and median \cite{Feuersaenger}.

Second important moments concerns data structure: the file should have a period as decimal separator (not a comma). Hence, because the initial Excel table had decimal points,  the comma were changed in the entire document to points using the simple method in Excel: “Find \& Replace”. Then the file was saved as .csv and read into \LaTeX \space according to the described code sample in the screenshots above. Third important technical moment concerns path to data. The .csv file should be saved in the same folder as the \LaTeX\space raw .tex file. Otherwise, the path should be declared explicitly. In terms of file management, it is also better to keep all the files of the current project in one directory.

\begin{figure}[H]
    \centering
    	\includegraphics[width=0.7\textwidth]{Fig-10.jpg}
    \caption{Flowchart of the research project: soils measurements. Source: authors.}
    \label{fig:10}
\end{figure}

Figure~\ref{fig:10} presents flowchart of the research project: soils measurements.

\section{Conclusion}

\subsection{Actuality}
Plotting 2D and 3D graphs accompanying calculations and computing tasks in engineering geology is an essential part of the work, since data visualization enables more effective evaluation of the correlations between the parameters and variables of soil samples. Graphical representation of the correlations of parameters of soils plotted against each other and compression strength allows to visually display physical and mechanical properties of soil samples. 

The importance of data visualization in engineering geology is clear: incorrect assessment of the soils properties may result to selecting wrong decision when constructing roads, railways and buildings. Negative environmental consequences may arise in case of the increased loads on the frozen soils in extreme climate which can lead to the crumbling road.

\subsection{Novelty}
Research novelty consists in technical approach of \TeX \space high-level, descriptive markup language application specifically for geological data processing and graphical visualization. Technical possibilities of engineering graphics by \TeX \space were demonstrated with screenshots of the codes used for plotting that can be applied in similar research. 

Correct decisions in constructing buildings and roads in engineering geology is largely base don the effective numerical methods that enable modeling soils compression, which is especially important for the construction works. Thus, plotting graphs and diagrams is an important part of the data analysis base don the laboratory data tables for the tasks of geological surveys. Plotting high-quality and accurate graphs is important assisting part in engineering and graphical works in geological analysis. 

Capabilities of \LaTeX \space demonstrated in this paper are by no means limited to the above examples. As a recommendation for future words it can be also noted that it is possible to visualize solid 3D surfaces by \LaTeX, as well as perform mathematical and statistical analysis of a larger amount of data in geotechnical surveys using special packages of \LaTeX.

\bibliographystyle{fcds_abbrv}

\begin{thebibliography}{2}
%
\bibitem{Gorokhovsky}
Gorokhovsky V. M., Tkachuk E. I., Modeling in engineering geology \emph{Novocherkassk, ed. NPI} - 1980. 84 pp.
%
\bibitem{Kaputin}
Kaputin Yu. E., Yezhov A. I., Henley S., \emph{Geostatistics in mining and geological practice}. Apatity, 1995.
%
\bibitem{Porotov}
Porotov G. S., \emph{Mathematical modeling methods in geology: Textbook}. SPbGGU (TU). St. Petersburg, 2006, 223 pp.
%
\bibitem{Feuersaenger}
Feuersaenger C., \emph{Manual for Package pgfplots 2D/3D Plots in \LaTeX}, Version 1.15.

\end{thebibliography}
\noindent{\emph{Please leave some space here for submission and acceptance date}}
\end{document}


